\documentclass[10pt]{article}
\usepackage[a4paper,margin=22mm]{geometry}
\usepackage[T1]{fontenc}
\usepackage[utf8]{inputenc}
\usepackage{lmodern,microtype}
\usepackage{amsmath,amssymb,bm,booktabs,array,longtable,graphicx}
\usepackage{natbib,caption,subcaption,xcolor,hyperref,placeins}
\hypersetup{colorlinks=true,linkcolor=blue!55!black,citecolor=blue!55!black,urlcolor=blue!55!black}
\setlength{\emergencystretch}{2em}
\setlength{\parskip}{0.25em}
\setlength{\parindent}{1em}
\renewcommand{\topfraction}{.92}
\renewcommand{\bottomfraction}{.85}
\renewcommand{\textfraction}{.07}
\renewcommand{\floatpagefraction}{.65}
\captionsetup{font=small,labelfont=bf}
\newcommand{\kms}{\mathrm{km\,s^{-1}}}
\newcommand{\Msol}{\mathrm{M}_{\odot}}
\newcommand{\E}{\mathbb{E}}
\newcommand{\ind}{\mathbf{1}}
\newcommand{\figfile}[2]{\IfFileExists{figures/#1.pdf}{\includegraphics[width=#2\linewidth]{figures/#1.pdf}}{\fbox{\parbox[c][35mm][c]{.88\linewidth}{\centering Figure #1}}}}
\newcommand{\tallfigfile}[1]{\IfFileExists{figures/#1.pdf}{\includegraphics[width=\linewidth,height=.72\textheight,keepaspectratio]{figures/#1.pdf}}{\figfile{#1}{.94}}}
% Edit authorship here before submission. Do not infer affiliations from a repository.
\newcommand{\PaperAuthors}{Nalin Dhiman}
\newcommand{\PaperAffiliation}{}


\makeatletter
\renewcommand*\l@section{\@dottedtocline{1}{0em}{2.8em}}
\renewcommand*\l@subsection{\@dottedtocline{2}{1.5em}{3.8em}}
\makeatother
\renewcommand{\thesection}{S\arabic{section}}
\renewcommand{\theequation}{S\arabic{equation}}
\renewcommand{\thefigure}{S\arabic{figure}}
\renewcommand{\thetable}{S\arabic{table}}
\title{Supplementary Material\\Measuring spiral structure and testing its incremental information for galaxy redshift estimation}
\author{\PaperAuthors\PaperAffiliation}
\date{}
\begin{document}
\maketitle
\begin{abstract}
This supplement provides the measurement definitions, reference conventions,
learning design and numerical details needed to interpret and reproduce the
main results. It also develops the physical boundary between image-based
redshift prediction and dynamical distance inference. The finite-equilibrium
and Tully--Fisher calculations are supporting analyses with independent
kinematic inputs; they are not features of the photometric-redshift model.
\end{abstract}
\tableofcontents

\section{Samples, identities and analysis units}
The principal annotation source is Galaxy Zoo:3D \citep{masters2021}, combined
with native Legacy Survey DR9 science images \citep{dey2019}. Observed catalog
features and the inherited heliocentric \texttt{NSA\_Z} label come from official
MaNGA DRPall within SDSS-IV DR17 \citep{york2000,bundy2015,dr172022}.
Identity joins precede feature assembly, so incidental image-download or row
order cannot assign one galaxy's features to another galaxy's target.
No physical radius, catalog stellar mass, mass inferred from a rotation curve
or reference distance enters the redshift design matrix.

The initial strong-spiral cohort has 160 objects with a spiral-vote fraction
of at least 0.7 and the stipulated single-center selection. It excludes 32
objects used for earlier feasibility work. Tile membership is
\begin{equation}
 g(\alpha,\delta)=\left(\left\lfloor\frac{\alpha}{10^{\circ}}\right\rfloor,
 \left\lfloor\frac{\delta+90^{\circ}}{10^{\circ}}\right\rfloor\right),
\end{equation}
with right ascension on its standard wrap. Each tile belongs to one split.
The catalog split has 91/40/29 galaxies and 38/13/13 sky groups for
training/calibration/test. The separate additional acquisition excludes all
192 earlier identities and all 67 earlier tiles, leaving 99 objects in 39
new groups. All 59 eligible strong spirals in that acquisition form the
principal additional cohort, in 30 groups. The other 40 objects describe
weaker or uncertain morphology and are not verified arm-free negatives.

The units of analysis differ by question. For region AP the basic outcome is
one score per galaxy. For the primary redshift contrast it is a mean paired
effect within each sky group. For centerline validation it is an entire galaxy
with its curve sets. None treats correlated pixels as independent galaxies.
Hash-based selection and spatial blocking make a pilot reproducible; they do
not yield a volume-complete or survey-representative sample. Blocking is a
general response to dependence \citep{roberts2017}, whose residual strength
must still be examined in a future larger study.

\begin{table}[htb]
\centering
\caption{Coverage denominators. ``Defined AP'' requires valid scored pixels
of both reference classes. The additional redshift analysis retains the
missing-image object through training-fitted imputation.}
\begin{tabular}{lrrr}
\toprule
Cohort & Objects & Sky groups & Defined arm AP\\
\midrule
Original training & 91 & 38 & training task\\
Original calibration & 40 & 13 & calibration task\\
Original test, three votes & 29 & 13 & 29\\
Additional strong spirals, three votes & 59 & 30 & 57\\
Additional strong spirals, five votes & 59 & 30 & 53\\
Separate geometry development/test & 12/12 & not the same split & centerline task\\
\bottomrule
\end{tabular}
\end{table}

\section{Coordinates, masks and photometric conventions}
\subsection{Native sampling and annotation projection}
The science grid is the native common $r$-band Legacy grid, approximately
0.262 arcsec per pixel. Science and inverse-variance planes must have matching
band, brick, dimensions and celestial WCS; all three bands must share the
registered grid and contain the requested galaxy center. Coordinate consistency
is checked on a $3\times3$ grid including corners, edges and center. Nonfinite
coordinates or an excessive interband displacement cause an explicit failure.
An image extension is selected from audited metadata, rather than by assuming
the last FITS extension contains the target.

The Galaxy Zoo:3D reference masks retain their original annotation coordinate
system. For each native science pixel, WCS transforms its sky coordinate to
the annotation grid and nearest-neighbor sampling transfers the vote count.
The science flux is not resampled in this primary pipeline. Nearest-neighbor
projection is suitable for preserving discrete label semantics, but does not
create information at the finer native sampling. Comparisons remain limited
by the reference resolution.

The allowed science region is a circular 3--23 arcsec annulus about the supplied
human center. Any positive foreground-star reference masks that pixel in every
band. Invalid inverse variance, nonfinite pixels and array boundaries are
also excluded. One common feature mask is used across every band and scale.
Reference coverage additionally limits \emph{scoring} and supervised pixel
sampling; it does not limit the image-defined annulus used to form downstream
redshift summaries.

\subsection{Background and radial structure}
A border of width $\max[2,\lfloor0.12\min(N_x,N_y)\rfloor]$ pixels supplies
a robust background estimate. Up to five iterations remove values farther
than 3.5 robust standard deviations from the median. The robust scale is
\begin{equation}
 \widehat\sigma_{\rm bg}=1.4826\,
 \operatorname{median}|I-\operatorname{median}(I)|.
\end{equation}
This statistic measures local image scatter, not a complete spatial noise
covariance. Extended galaxy light can contaminate a small cutout's border.

The radial profile is the median in one-pixel-wide radial bins, using valid
pixels and requiring at least four values in a bin. Missing bins are
interpolated and the profile is smoothed with a one-bin Gaussian. Subtracting
this profile yields the residual $J$. A nearest-valid fill supplies values
needed to evaluate convolutions internally; filled pixels and their guarded
neighborhoods never become accepted feature pixels. For the original supervised
region experiment the radial coordinate is projected and circular. An elliptical
profile is used only when an independently specified disk geometry is supplied
for the distinct geometric analysis.

Catalog fluxes in nanomaggies become magnitudes via
\begin{equation}
 m_b=22.5-2.5\log_{10}F_b,\qquad
 \sigma_{m_b}\simeq\frac{2.5}{\ln10}\frac{\sigma_{F_b}}{F_b}.
\end{equation}
Nonpositive or unavailable catalog fluxes become missing magnitude values,
handled by training-only imputation. Pixel color descriptors instead use
signed $\operatorname{asinh}$ differences, retaining negative noisy fluxes.
An annular ``positive flux'' is explicitly the sum of
$\max(I-B,0)$ over valid pixels; it is an image descriptor, not a claim of
unbiased total photometry.

\section{PSF matching and noise propagation}
Let the circular moment-equivalent Gaussian width in band $b$ be $w_b$ in
FWHM units. Matching to $w_* =\max_b w_b$ uses
\begin{equation}
 \sigma_{k,b}=\frac{\sqrt{w_*^2-w_b^2}}{2\sqrt{2\ln2}},
\end{equation}
in angular units. No sharpening is performed. For a normalized kernel $K$
and validity indicator $M$, masked convolution is
\begin{equation}
 \widetilde I_j=\frac{\sum_k K_{jk}M_kI_k}{a_j},\qquad
 a_j=\sum_k K_{jk}M_k.
\end{equation}
Assuming independent input noise, the propagated diagonal variance is
\begin{equation}
 \widetilde V_j=\frac{\sum_kK_{jk}^2M_kV_k}{a_j^2}.
 \label{eq:variance}
\end{equation}
This requires the \emph{squared} kernel. It does not make the smoothed pixels
independent. Their covariance is
\begin{equation}
 \operatorname{Cov}(\widetilde I_j,\widetilde I_l)
 =\frac{\sum_k K_{jk}K_{lk}M_kV_k}{a_ja_l}
\end{equation}
under the same idealization. Native coadds may already have covariance not
represented in the inverse-variance plane.

Pixels with convolution support below 0.99 are excluded. After this matching,
the feature mask also requires more than four times the largest feature
$\sigma$ from holes and borders: the 1.6 arcsec filter therefore requires a
6.4 arcsec guard. At least 100 common valid feature pixels are required.
Coverage is measured after these guards, rather than silently filling masked
areas with model predictions. Bands are not coadded for the primary classifier;
each band retains its own feature block.

A deterministic diagnostic using the first development galaxy illustrates
the difference between moment matching and complete PSF matching. The three
matched moment-equivalent FWHM values agree to approximately $10^{-5}$ in
relative spread, but their positive-stamp half-light radii still span
0.636--0.735 arcsec and their 80\%-light radii 1.139--1.299 arcsec. Approximate
matching therefore controls one resolution statistic without equalizing the
core or wings. This one-object diagnostic is a method check, not a population
error bound or a basis for retuning the frozen test analysis.

\section{Exact local feature representation}
Write $N=I-B$ for the background-subtracted native signal, $J=I-m(R)$ for the
radial residual, and $S(X)=\max[1.4826\operatorname{MAD}(X),0.1
\operatorname{sd}(X),10^{-12}]$ for the stabilizing scale on common feature
pixels. Let $s_N=S(N)$, $s_J=S(J)$ and
$f=\max(\widehat\sigma_{\rm bg},0.1s_N,10^{-12})$. The brightness median
is $\mu_N$. The filter standard deviations are
$\sigma_1=0.4$, $\sigma_2=0.8$, $\sigma_3=1.6$ arcsec, converted to pixels
before discrete convolution. Derivative factors use pixel-coordinate
$\sigma$ consistently.

For each band let $L_k=G_{\sigma_k}*J$ and
\begin{equation}
 H=\sigma_2^2\begin{pmatrix}
 \partial_{xx}L_2&\partial_{xy}L_2\\
 \partial_{xy}L_2&\partial_{yy}L_2
 \end{pmatrix}.
\end{equation}
Order its eigenvalues by absolute magnitude, $|\lambda_s|\leq|\lambda_l|$.
The bright-ridge strength is
\begin{equation}
 \mathcal R=\max(-\lambda_l,0)
 \exp\left[-\frac12\left(\frac{\lambda_s}
 {0.5\max(|\lambda_l|,10^{-12})}\right)^2\right].
\end{equation}
It is a local descriptor inspired by established Hessian approaches
\citep{frangi1998}, not a new physical definition of an arm.
Scale-normalized filtering is related to the broader scale-space framework
of \citet{lindeberg1998}; this experiment uses three fixed scales rather than
automatic physical-scale selection.

\begin{table}[htb]
\centering\small
\caption{Thirteen features repeated in each of $g,r,z$. Together with three
radius terms and two color terms they give 44 features. All are computed
without redshift or physical-distance information.}
\begin{tabular}{p{.32\linewidth}p{.59\linewidth}}
\toprule
Feature & Definition or meaning\\
\midrule
Signed intensity & $\operatorname{asinh}(N/f)$\\
Normalized brightness & $(N-\mu_N)/s_N$\\
Radial contrast & $J/s_J$\\
Radial profile & $(N-J)/s_N$\\
Fine/middle contrast difference & $(L_1-L_2)/s_J$\\
Middle/coarse contrast difference & $(L_2-L_3)/s_J$\\
Negative Laplacian-like response & $-\operatorname{tr}H/s_J$\\
Bright ridge & $\mathcal R/s_J$\\
Gradient strength & $\sigma_1|\nabla L_1|/s_J$\\
Fine residual & $L_1/s_J$\\
Coarse residual & $L_3/s_J$\\
Radial orientation, even & $(\mathcal R/s_J)\cos 2(\alpha-\phi_{\rm sky})$\\
Radial orientation, odd & $(\mathcal R/s_J)\sin 2(\alpha-\phi_{\rm sky})$\\
\bottomrule
\end{tabular}
\end{table}

Here $\alpha$ denotes the local bright-ridge tangent orientation and
$\phi_{\rm sky}$ the projected center-to-pixel angle. Double angles remove the
arbitrary sign of an eigenvector; these features do not measure a globally
deprojected spiral pitch. The three radius terms are $u$, $u^2$ and
$\log(1+u)$, with $u=R_\theta/23$ arcsec. The two color terms are
\begin{equation}
 C_{ab}=\operatorname{asinh}(N_a/f_{ab})-
 \operatorname{asinh}(N_b/f_{ab}),\qquad
 f_{ab}=\max(\widehat\sigma_a,\widehat\sigma_b,10^{-12}),
\end{equation}
for $(a,b)=(g,r),(r,z)$ in common declared photometric units. This signed
color is useful for noisy pixels but is not a standard magnitude difference.

\section{Reference classes, learner candidates and calibration}
The primary binary label is $y_j=\ind[v_j\geq3]$, where $v_j$ is the number
of human arm polygons covering the scored pixel. The classifier is trained
against this operational definition. Because the reference is incomplete,
$y=0$ denotes the comparison class under the task, not verified physical
absence. The selection already favors strong spirals; no real-negative
specificity is inferred from low-vote or unasked objects.

Each training galaxy supplies the same number of randomly selected positive
and comparison pixels, $\min(400,n_+,n_-)$. Sampling uses a deterministic
identity-derived seed and inverse sample-size weights give each usable galaxy
the same total weight. Candidate evaluation averages per-galaxy AP within
three-fold training sky-group cross-validation. The seven candidates are:
brightness/radius logistic regression; brightness/radius ExtraTrees at depth
8 or 12; full logistic regression at $C=0.1$ or $1$; and full ExtraTrees at
depth 8 or 12. Standardization for logistic regression is fitted within the
training fold. Tree models use 80 trees, leaf size 32, feature fraction 0.7
and a fixed seed. The randomized tree construction follows
\citet{geurts2006} through scikit-learn \citep{pedregosa2011}.

The selected full and brightness/radius models both have depth 12. Their
training-CV AP values are 0.6757 and 0.6335. Thus the original test comparison
is between independently training-selected models with a common eligible
capacity family, not between a complex learner and an intentionally weak
baseline. Pixel-class balancing does change score prevalence and motivates
distinguishing raw ranking scores from calibrated annotation probabilities.

The display mapping takes raw score $s$ to
\begin{equation}
 p_{\rm display}=\operatorname{logistic}
 [a\operatorname{logit}(\operatorname{clip}(s))+b],\qquad a>0.
\end{equation}
This is a sigmoid calibration \citep{platt1999}; the fitted values are
$a=0.81289$, $b=-1.48564$. The calibration-selected display threshold is 0.30.
The positive slope preserves AP ranking. A different artifact retains raw
$s$ and fixed threshold 0.5 for downstream inference, so these calibration
labels cannot influence the redshift representation.

For ranked pixels with recall increments $R_k-R_{k-1}$ and precision $P_k$,
\begin{equation}
 \operatorname{AP}=\sum_k(R_k-R_{k-1})P_k.
\end{equation}
AP depends on reference prevalence as well as ranking. A threshold change
from three to five polygons changes prevalence and occasionally eliminates
a scoring class; it is not a fixed-task perturbation alone. On the original
test, five-vote AP is 0.5012 on 28 objects; on the additional strong cohort it
is 0.4868 on 53 objects. The corresponding primary three-vote AP values are
0.7197 on 29 and 0.6890 on 57. These denominators are part of the result.

\begin{figure}[tb]
\centering\figfile{arm_consensus_sensitivity}{.94}
\caption{Sensitivity to the human consensus threshold. Three- and five-vote
targets have different positive prevalence and evaluable samples, so changes
in AP cannot be interpreted as a controlled loss of physical arm accuracy.
Low-vote/unasked regions are not promoted to verified negatives.}
\end{figure}

\section{Downstream features and cross-fitting}
Training arm summaries are out-of-fold by sky group. If $h^{(-k)}$ is the
classifier fitted without fold $k$ and $\mathcal A$ is the image-defined
annulus, a training galaxy in that fold receives
\begin{equation}
 A_i=\mathcal T\{h^{(-k)}(X_{ij}):j\in\mathcal A_i\}.
\end{equation}
Calibration and test galaxies instead receive $\mathcal T$ applied to the
single full-training classifier. This is cross-fitting of the arm
representation. Upstream candidate choice uses the training set, so downstream
training-CV results conditional on these cached descriptors are not a fully
nested assessment of the complete model-selection process. The original
final test groups remain outside all training selection.

For raw scores $s_j$ and normalized radii $u_j$, define
\begin{align}
 \bar s&=n^{-1}\sum_j s_j,&
 u_s&=\frac{\sum_j s_ju_j}{\sum_j s_j},\\
 H_s&=-n^{-1}\sum_j[s_j\log_2s_j+(1-s_j)\log_2(1-s_j)],&
 f_s&=n^{-1}\sum_j\ind[s_j\geq0.5].
\end{align}
Together with $Q_{0.9}(s)$, $\operatorname{sd}(s)$ and inner/outer score means
at $u=0.55$, these give eight summaries. Zero-weight and empty subregion
conventions are fixed by the implementation. Entropy here is a feature of
the score distribution; it is not a posterior uncertainty on physical arms.

The four regressor families are nested in their principal information blocks:
P, PG, PGI and PGIA. P has five magnitudes and adjacent colors. PG adds
angular half-light radius and axis ratio. PGI adds three annular positive-flux
sums, the $r$-band half-positive-flux radius, $r$ residual contrast, and two
mean signed colors. PGIA adds $A_i$. All share five magnitude errors, three
PSF widths, three native noise estimates and an arm-feature availability flag.
Training-median imputation, including missingness indicators, is fitted
within training. No test value determines an imputation median.

The target transformation $t=\log(1+z)$ preserves order and compresses the
redshift range. It does not cause the regressor to learn a cosmological
distance. Predictions $\hat z=\exp(\hat t)-1$ are assessed against inherited
spectroscopic catalog labels. Redundant magnitudes and colors are intentionally
kept fixed across ablations; the goal is a controlled contrast, not a claim
that this small ExtraTrees system is an optimal photometric-redshift architecture.
Image-based CNN inference and larger morphology-assisted methods provide
relevant alternatives \citep{pasquet2019,soo2018}.

\section{Group-conformal interval construction}
Let $\widehat t_i$ be the prediction from a model fitted only to training
galaxies. Each calibration group contributes
\begin{equation}
 S_g=\max_{i\in g}|t_i-\widehat t_i|.
\end{equation}
For $G_c$ calibration groups and miscoverage $\alpha$, define
\begin{equation}
 k=\left\lceil(G_c+1)(1-\alpha)\right\rceil,
 \qquad q=S_{(k)},
\end{equation}
where $S_{(k)}$ is the $k$th ordered group score; a rank beyond $G_c$ requires
an infinite interval. Here $G_c=13$, $\alpha=0.1$ and $k=13$. Prediction
intervals are transformed monotonically:
\begin{equation}
 [z_-,z_+]=[\max(0,\exp(\widehat t-q)-1),
 \exp(\widehat t+q)-1].
\end{equation}
Under exchangeability of the calibration and future group objects, the rank
argument supplies coverage for the group maximum, hence simultaneous coverage
of all retained members of a new comparable group. It does not give precise
conditional coverage at each feature value. The basic split-conformal theory
is described by \citet{lei2018}; \citet{angelopoulos2021} gives a practical
account. Our calculation does not implement the likelihood-ratio correction
that can address a specified covariate shift \citep{tibshirani2019}.

The calibrated set is small, so the maximum residual can produce broad
intervals. Original P coverage 29/29 and additional P coverage 57/59, with
median widths 0.08954 and 0.09043, are empirical diagnostics, not evidence of
precise physical distances. The additional cohort has whole-group coverage
28/30. Its selection and missingness can differ from the calibration groups;
coverage is therefore measured rather than assumed to transfer.

\section{Effect sizes, estimands and prospective planning}
The primary paired effect averages absolute-error differences within each
sky group and then gives groups equal weight. If group sizes differ,
\begin{equation}
 \frac{1}{G}\sum_g\frac{1}{n_g}\sum_{i\in g}d_i
 \ne \frac{1}{N}\sum_i d_i
\end{equation}
in general. The same distinction applies to paired AP improvements. We report
galaxy-macro measurements for descriptive performance and equal-group paired
effects for the central comparison. Percentile intervals use whole-group
bootstrap resampling with 100,000 draws for the final redshift contrasts
\citep{efron1979}. They condition on the fitted models and cohorts.

\begin{table}[htb]
\centering\small
\caption{Principal equal-group redshift contrast. Positive values favor arms.
The Student-$t$ interval is a sensitivity to small group counts, not an
independent replication.}
\begin{tabular}{lrrr}
\toprule
Cohort & Effect & Bootstrap 95\% & Student-$t$ 95\%\\
\midrule
Original, $G=13$ & $-0.000471$ & $[-0.001391,0.000367]$ & $[-0.001491,0.000549]$\\
Additional, $G=30$ & $0.000440$ & $[-0.000555,0.001452]$ & $[-0.000620,0.001500]$\\
\bottomrule
\end{tabular}
\end{table}

The practical margin $\delta_z=0.001$ was introduced retrospectively. A second
sensitivity uses $\delta_t=0.00138663$ in $t=\log(1+z)$, equal to 10\% of a
specified training-only PGI CV error on that transformed scale. It must not
be mislabeled as a redshift margin. Equivalence testing reverses the ordinary
null: it seeks to exclude effects outside a prespecified interval, rather than
merely failing to reject zero \citep{schuirmann1987}. No primary useful-gain
or equivalence conclusion is established here. The secondary transformed-scale
margin is more permissive and remains a sensitivity, not a newly confirmatory
success.

For a frozen future comparison with independent approximately normal group
effects of standard deviation $s_d$, the usual planning scale for a two-sided
test against zero is
\begin{equation}
 G\ \approx\ \frac{(z_{1-\alpha/2}+z_{1-\beta})^2s_d^2}{\delta_z^2}.
\end{equation}
The numerical planning uses a noncentral-$t$ calculation instead of this
large-sample approximation. With $s_d=0.00283859$, the larger observed cohort
SD, detecting a true gain 0.001 requires 66 groups at 80\% power and 87 at
90\%. With the normal-assumption upper-variance scenario
$s_d=0.00363255$, these become 106 and 141 groups. A test of superiority over
0.001 when the true gain is 0.002, or an equivalence design around zero, is
a different calculation. These figures are planning assumptions, not estimates
of observed power or predictions that morphology will help.

\subsection{Literature-motivated observable controls}
The L5 control supplements P and apparent geometry with observed $r$-band
surface brightness and catalog S\'ersic index. Surface brightness is
\begin{equation}
 \mu_r=m_r+2.5\log_{10}(2\pi q R_{50}^2),
\end{equation}
with $R_{50}$ in arcsec and apparent axis ratio $q$. This observed quantity
uses no redshift-based dimming or rest-frame correction. L3 restricts explicit
photometric and magnitude-error predictors to $grz$, with the stated apparent
geometry and $grz$ quality information retained. Appending A gives L5A/L3A.
These are common-capacity controls motivated by \citet{soo2018}; they do not
reproduce that paper's neural architecture or CS82 sample.

\begin{table}[htb]
\centering\small
\caption{Retrospective controls on the same 59 additional strong spirals.
All are fitted only on the original 91 training galaxies. Group and galaxy
estimands can change in different directions, as for L3/L3A.}
\begin{tabular}{lrr}
\toprule
Family & Galaxy MAE & Equal-group MAE\\
\midrule
L5 & 0.012463 & 0.013132\\
L5A & 0.012393 & 0.012547\\
L3 & 0.012504 & 0.012928\\
L3A & 0.012634 & 0.012786\\
\bottomrule
\end{tabular}
\end{table}
The equal-group L5 gain is $0.000584$ with 95\% interval
$[-0.000331,0.001658]$; the L3 gain is $0.000142$ with interval
$[-0.000802,0.001195]$. No test-based deployed-model reselection or new
interval-calibration claim is made for these additional baselines.

\section{Centerline validation and reference uncertainty}
The geometry sample contains 24 Galaxy Builder subjects with human freehand
traces, author-cleaned clustered arms and fitted disk geometry
\citep{lingard2020,lingard2021}. The SDSS $r$ image arrays have normalized
linear values, PSFs and 0.396 arcsec per pixel metadata. They are not
absolutely calibrated three-band flux images. Source image centers exclude
earlier objects within 10 arcsec, but the geometry sample is not a new
survey or a tile-disjoint extension of the principal experiment.

Twelve development objects choose one smoothing/template configuration.
The extractor searches logarithmic-spiral contrast templates after
deprojection with the supplied geometry. Twelve test objects evaluate the
fixed configuration. Source references were available during preparation;
there was no sealed external test. Scores were not used to tune the
configuration after the test evaluation. This distinction matters when
planning a future confirmatory geometry benchmark.

Let $P$ and $T$ be sampled predicted and reference centerline point sets on
common annular support, and let $d(p,T)=\min_{t\in T}|p-t|$. At angular
tolerance $\epsilon$ the point-based metrics are
\begin{equation}
 \operatorname{Prec}_\epsilon=\frac{1}{|P|}\sum_{p\in P}\ind[d(p,T)\leq\epsilon],
 \quad
 \operatorname{Rec}_\epsilon=\frac{1}{|T|}\sum_{t\in T}\ind[d(t,P)\leq\epsilon],
\end{equation}
with $F1=2\operatorname{Prec}\operatorname{Rec}/
(\operatorname{Prec}+\operatorname{Rec})$. Curves are sampled approximately
uniformly in projected angular length at 0.1 arcsec spacing, with at least
eight points per curve. The main tolerance is the supplied PSF FWHM; half
that tolerance is a sensitivity. One-to-one segment assignment minimizes
the symmetric median point-to-curve distance. Only pairs with assigned
distance at most twice the PSF FWHM enter pitch comparison, preserving the
matched denominator.
Test F1 is 0.223 (95\% whole-galaxy interval 0.105--0.353); at half-PSF it
is 0.135 (0.049--0.235). There are 70 predicted and 19 reference segments,
with seven matched segments whose median absolute pitch difference is
11.42 degrees. A conditional pitch statistic cannot describe every unmatched
arm or galaxy.

Twenty of 24 synthetic nonspiral images produce at least one false curve.
This is a controlled stress test of the template behavior, not a measured
false-positive rate for a real observed nonspiral population. Human traces
also have uncertainty: different volunteers can select different arm pieces,
and author aggregation/fitting imposes a specific reference convention.
Independent tracing is stronger than visual plausibility, but not error-free
physical truth. No observed arm-width accuracy is established.

\section{Archival published-pitch comparison}
The \citet{hart2017} author catalog contains 6,222 parent galaxies, with
SpArcFiRe-based pitches after reliable-arc selection. Its exact DR7 identifiers
match all 24 Galaxy Builder subjects, with maximum coordinate separation
0.020 arcsec. Thirteen pitches are finite and eleven retain the published
$-999$ sentinel as missing. There are seven measured development and six
measured test objects. No missing pitch is scored as zero or as absence of arms.

The Galaxy Builder reference summary is an unsigned pitch averaged over
author-accepted arm fits using their stored \emph{projected native-image}
lengths. The archival SpArcFiRe summary uses selected reliable arcs with a
different standardized/deprojected weighting convention and chirality
selection. Equal galaxy summaries need not imply agreement about individual
arm locations or completeness. The arithmetic therefore measures a
convention-dependent catalog difference.

For the six measured test identities the mean absolute difference is
7.727 degrees, median absolute difference 7.280 degrees, and the conditional
whole-galaxy bootstrap interval for mean absolute difference is
5.061--10.685 degrees. On those same identities, a fixed development-median
constant has mean absolute difference 7.122 degrees. Seven development
identities give mean absolute difference 10.357 degrees; their signed mean
is 6.925 degrees. These different statistics are not interchangeable.

This is not an execution of the author runtime on the present images.
The identities share the Hart parent population and individual classifier
training membership is not established. The comparison is retrospective and
does not create an independent population test. It neither rescues the
failed centerline prototype nor establishes superiority over SpArcFiRe
\citep{davis2014,hart2017}. A future algorithm comparison must harmonize
input images, support, arc-selection definitions, failure handling and
training-overlap exclusions.

\section{Physical sizes, arm widths and error propagation}
For supplied external $D_A$, angular separations convert through
$L=D_A\theta$. First-order uncertainty, including a possible covariance, is
\begin{equation}
 \operatorname{Var}(L)\simeq\theta^2\operatorname{Var}(D_A)
 +D_A^2\operatorname{Var}(\theta)
 +2D_A\theta\operatorname{Cov}(D_A,\theta).
\end{equation}
A physical size reported without the supplied distance uncertainty can appear
more informative than it is. Disk deprojection adds uncertainty in center,
inclination and position angle. For example, $Y=y'/\cos i$ gives
$\partial Y/\partial i=y'\sec i\tan i$, exposing rapidly increasing
sensitivity as the disk approaches edge-on. A projected separation between
two galaxies is also not a three-dimensional spatial separation without
line-of-sight information.

For approximately Gaussian transverse profiles one sometimes writes
$w_{\rm obs}^2\simeq w_{\rm int}^2+w_{\rm PSF}^2$ in a common width
convention. Deconvolution near $w_{\rm obs}=w_{\rm PSF}$ is unstable and
requires uncertainty propagation, background control, a defensible local
cross-arm direction and a PSF model. Real arm profiles and PSFs need not be
Gaussian. Conditional fits around supplied reference curves can be useful
diagnostics, but they are not independent observed-width validation and are
not accepted distance predictors in this study.

\section{A globally finite dynamical equilibrium}
\label{sec:finite}
This supporting calculation asks whether the dynamical component can be
formulated consistently. It uses known equilibrium physics, not an additional
image-derived distance calibration. Pressure support in a stationary,
nonrotating isotropic halo satisfies
\begin{equation}
 \nabla P=-\rho\nabla\Phi,\qquad
 \nabla^2\Phi_h=4\pi G\rho,\qquad
 \Phi=\Phi_h+\Phi_b.
 \label{eq:hydro}
\end{equation}
An untruncated isothermal relation $P=\rho s^2$ gives
$\rho\propto\exp(-\Phi/s^2)$ but does not by itself specify a globally finite
isolated halo. Cutting such a profile at an arbitrary nonzero-pressure edge
requires a boundary stress. We instead implement a lowered-isothermal
closure with a finite zero-pressure surface, following the established
King family \citep{king1966,gieles2015}. Self-consistent disk/halo construction
also has a long history \citep{kuijken1995}.

\subsection{Distribution function and pressure moment}
Let $E=v^2/2+\Phi$, choose a cutoff energy $E_t$, and define
\begin{equation}
 f(E)=\begin{cases}
 A\{\exp[(E_t-E)/s^2]-1\},&E<E_t,\\
 0,&E\geq E_t.
 \end{cases}
 \label{eq:king}
\end{equation}
An ergodic function of conserved energy is a stationary collisionless
distribution in the prescribed static potential. Its velocity moments also
give an isotropic Jeans/barotropic equilibrium. Neither interpretation
requires a collisional viscosity. Set $W=(E_t-\Phi)/s^2$ and define, for $W>0$,
\begin{align}
 D(W)&=e^W\operatorname{erf}\sqrt W-
 2\sqrt{W/\pi}\left(1+\frac{2W}{3}\right),\\
 Q(W)&=e^W\operatorname{erf}\sqrt W-
 2\sqrt{W/\pi}\left(1+\frac{2W}{3}+\frac{4W^2}{15}\right).
\end{align}
Both are zero for $W\leq0$. With central density $\rho_0$ and central binding
$W_0$, the moments are
\begin{equation}
 \rho(W)=\rho_0\frac{D(W)}{D(W_0)},\qquad
 P(W)=\rho_0s^2\frac{Q(W)}{D(W_0)}.
\end{equation}
Since $Q'(W)=D(W)$ and $\nabla W=-\nabla\Phi/s^2$, equation~\eqref{eq:hydro}
follows directly. The local one-dimensional dispersion is $P/\rho$, not
the fixed parameter $s^2$. Near the edge,
\begin{equation}
 D(W)\sim\frac{8W^{5/2}}{15\sqrt\pi},\qquad
 Q(W)\sim\frac{16W^{7/2}}{105\sqrt\pi},\qquad
 \frac{P}{\rho}\sim\frac{2}{7}s^2W.
\end{equation}
Density and pressure therefore vanish continuously on the cutoff
equipotential. The cutoff is a specified model family, not an inference
of a cosmological halo edge or formation history.

\subsection{Flattened baryons and isolated gravity}
The prescribed baryonic potential is a nonnegative superposition of
Miyamoto--Nagai components \citep{miyamoto1975},
\begin{equation}
 \Phi_b(R,z)=-\sum_j\frac{GM_j}
 {\sqrt{R^2+[a_j+\sqrt{z^2+b_j^2}]^2}},\qquad M_j\geq0.
\end{equation}
Its equatorial circular-speed contribution is
\begin{equation}
 v_{b,j}^2(R)=\frac{GM_jR^2}{[R^2+(a_j+b_j)^2]^{3/2}}.
\end{equation}
Thus equatorial forces determine $a_j+b_j$ more directly than the vertical
thickness $b_j$. The halo density is iterated in the combined potential, so
adding a disk force after solving an isolated spherical halo is not mislabeled
as exact coupled equilibrium.

For axisymmetric density expanded as $\rho(r,\mu)=\sum_l\rho_l(r)P_l(\mu)$,
the isolated Green solution is
\begin{equation}
 \Phi_{h,l}(r)=-\frac{4\pi G}{2l+1}
 \left[r^{-l-1}\int_0^r\rho_l(a)a^{l+2}\,da
 +r^l\int_r^\infty\rho_l(a)a^{1-l}\,da\right].
\end{equation}
Reflection symmetry retains even $l$. Quadrature and under-relaxed density
iteration close Poisson's equation with the moments above. The numerical
sphere must enclose the entire zero-density surface; a state with positive
boundary binding is rejected or the domain enlarged. Outside the halo the
potential continues through the vacuum multipoles, not an imposed zero-potential
wall. The analytic baryon components have finite mass but need not be
compactly supported inside the same sphere.

\begin{figure}[tb]
\centering\figfile{finite_halo_profiles}{.94}
\caption{Finite lowered-isothermal equilibrium and numerical closure.
The cutoff surface has zero density and pressure; the computational boundary
lies in vacuum. This controlled calculation concerns the specified equilibrium
family and baryonic potential, not identification of dark-matter microphysics.}
\end{figure}

\subsection{Numerical validation and three-galaxy application}
An independent spherical ordinary-differential-equation solution checks the
zero-baryon limit. The controlled case has maximum relative circular-speed
difference $3.55\times10^{-5}$ and total-mass difference
$5.36\times10^{-5}$. The direct pressure-gradient identity has maximum
relative residual $1.77\times10^{-9}$. Increasing the controlled grid from
200 radial points/$l_{\max}=8$ to 960/$l_{\max}=32$ reduces the measured
Poisson residual from $7.44\times10^{-3}$ to $2.92\times10^{-4}$; the
coarse-to-finest speed difference is 0.0414 $\kms$. Doubling the vacuum
domain from 250 to 500 kpc changes speed by at most
0.000174 $\kms$ and mass by $5.43\times10^{-6}$ in relative units.
These are discretization checks for controlled inputs, not uncertainty
estimates for real galaxies.

For NGC2403, NGC3198 and NGC5055, SPARC supplies external distances,
photometric/gas force columns and observed rotation curves \citep{lelli2016}.
Nonnegative analytic components approximate the baryonic force columns without
using observed velocities to define the baryon template. This approximation
has baryonic speed errors of several $\kms$ and cannot reproduce every
outward-force feature of a ring-like gas distribution. Published signed gas
forces are preserved in the original-force comparison.

A photometric population relation sets
\begin{equation}
 h_z=10^{-0.708}(R_d/\mathrm{kpc})^{0.633}\,\mathrm{kpc},
\end{equation}
related to the disk-flattening information discussed by
\citet{bershady2010} and adopted in the SPARC structural context
\citep{lelli2016}. It gives 0.241, 0.404 and 0.409 kpc. The stellar analytic
model's first absolute vertical moment is matched to this value at $2.2R_d$.
This is a population prior, not a direct thickness measurement for each target.
A sensitivity varies the adopted height by $e^{\pm0.25}$ at fixed fitted dark
parameters; the largest resulting speed change is 0.494 $\kms$. The gas
thickness remains an assumed fraction, not a new observation.

The same predeclared image annuli define a regional conditional prediction
test: fit velocities outside the annulus and predict inside, using externally
supplied distance to convert angles to kpc. These annuli have no independent
arm-centerline acceptance and must not be called verified arms. They were
already examined in earlier analysis, so the comparison is exploratory.
The primary finite coupled model fixes $W_0=5$ and fits two dark parameters,
$\rho_0,s$. Table~\ref{tab:fluid} keeps the same three targets and all comparison
outcomes. No derivative of observed velocity is used as a predictor of that
same velocity.

\begin{table}[htb]
\centering\small
\caption{Regional held-out rotation-speed RMSE in $\kms$. All comparisons
condition on external distance and baryonic assumptions. Different baryonic
force approximations are explicitly distinguished; the three targets do not
establish population superiority.}
\label{tab:fluid}
\begin{tabular}{lrrr}
\toprule
Halo/force construction & NGC2403 & NGC3198 & NGC5055\\
\midrule
Finite coupled, $W_0=5$, analytic baryons & 4.69 & 11.77 & 18.13\\
Confined coupled isothermal, same thickness & 7.47 & 3.57 & 27.90\\
Finite spherical + analytic baryon forces & 8.32 & 18.01 & 9.05\\
Uncontracted isothermal + analytic forces & 5.61 & 5.58 & 10.44\\
Uncontracted isothermal + signed SPARC forces & 5.64 & 4.59 & 6.52\\
\bottomrule
\end{tabular}
\end{table}

Physical consistency does not guarantee a better conditional prediction. The
cutoff also matters: changing $W_0$ to three or seven at fixed dark parameters
can change speeds by tens of $\kms$. A later training-only cutoff search
over three, five, seven and nine reaches endpoints for these targets. This
does not identify a unique physical cutoff or establish an observational
microphysical inference.

\subsection{Distance scaling and why viscosity is not measured}
For free physical mass and length parameters, Newtonian equilibrium admits
the similarity
\begin{equation}
 r\mapsto\lambda r,\quad M\mapsto\lambda M,\quad
 \rho\mapsto\lambda^{-2}\rho,\quad
 \Phi\mapsto\Phi,\quad v_c\mapsto v_c.
\end{equation}
The same angular velocities can then correspond to differently scaled
systems. The implemented similarity check preserves speed to
$1.08\times10^{-14}$ in relative units. This statement concerns freely
rescaled mass models. If apparent flux is fixed and a calibrated
mass-to-light ratio is supplied, baryonic mass instead scales as $D^2$;
that extra physical calibration and independent kinematics can constrain
distance. The similarity must not be used to deny this conditional route.

For a rotating Newtonian fluid the cylindrical shear relevant to viscous
angular-momentum transport is
\begin{equation}
 \tau_{R\phi}=\eta R\frac{d\Omega}{dR}
 =\eta\left(\frac{dv_\phi}{dR}-\frac{v_\phi}{R}\right).
\end{equation}
Rigid rotation has zero shear; a flat rotation curve has nonzero shear even
though $dv_\phi/dR=0$. An isotropic pressure-supported halo and a gravitational
rotation curve do not identify $\eta$. A transport claim would require a
specified moving fluid, momentum/energy closure, time dependence or radial
flow and boundary torque information. The present finite model supplies
equilibrium density and gravity, not those missing observations.

\section{A calibrated kinematic distance route}
The Tully--Fisher relation adds independent kinematic information to photometry
\citep{tully1977}. For a baryonic version write
\begin{equation}
 y=\log_{10}(M_b/\Msol)=a x+b+\epsilon,\qquad
 x=\log_{10}(V/V_0),\quad
 \epsilon\sim\mathcal N(0,\sigma_{\rm int}^2).
\end{equation}
Characteristic velocity must retain its definition. Flat outer H\,I speed,
optical rotation speed and integrated line width are different observables,
as quantified by \citet{lelli2019}. Calibrated baryonic distances already
have an established literature \citep{schombert2020}; this supporting
implementation is not a new distance law.

Let $C_b=M_b/D^2$ be the distance-free mass/flux coefficient for a fixed
stellar mass-to-light convention and gas correction. Then
\begin{equation}
 \log_{10}(D/\mathrm{Mpc})=
 \frac12[a x+b-\log_{10}C_b],
\end{equation}
with unit conventions retained in $C_b$. Recovering $C_b$ from a published
flux-derived mass and its \emph{same reference distance} cancels that distance
algebraically; recovering it from inconsistent catalogs would not. Across
153 matched published rows, reconstructed
$M_b=0.5L_{3.6}+1.33M_{\rm HI}$ agrees with the rounded published masses to
at most 0.005 dex. This checks source consistency, not the correctness of
their distance anchors. A coherent test $D\mapsto2.5D$, $M_b\mapsto6.25M_b$
preserves $C_b$ to floating-point precision.

For Gaussian errors and a flat latent-log-speed approximation, integrating
over uncertain $x$ gives a vertical likelihood with variance
\begin{equation}
 s_i^2=\sigma_{y,i}^2+a^2\sigma_{x,i}^2+\sigma_{\rm int}^2.
\end{equation}
The implementation uses bounded slope and scatter priors and propagates
posterior and target flux/speed uncertainty. It does not fit a population
selection function or shared anchor covariance. Distance-related mass error
and reference-distance error must not be independently counted twice when
constructing the flux coefficient.

Thirty calibrators with external stellar-distance reference methods are
evaluated leave-one-out. The
error-aware model has median absolute fractional distance error 17.87\%,
compared with 17.60\% for ordinary least squares on the same updated data.
Its nominal 90\% model intervals cover 27/30 reference points, with median
width 82.21\% of the predicted distance. This is a useful calibrated route,
but the broad intervals and small reused anchor sample preclude a claim of
new precision or externally calibrated coverage.

\begin{figure}[tb]
\centering\figfile{btfr_validation}{.94}
\caption{Supporting baryonic Tully--Fisher distance calculation. Distances
require matched flux conventions, independent measured kinematics and calibrated
anchors. The plotted intervals are model-dependent; the small local
leave-one-out sample is not a new blind distance validation.}
\end{figure}

A retrospective matched-width audit keeps the same 23 target galaxies and
22 leave-one-out calibrators for all alternatives. Proper line-width
calibration gives median absolute error 23.13\%, mean absolute error 24.60\%
and interval coverage 22/23. Deliberately substituting that width into a
flat-speed calibration happens to lower median error to 14.82\%, but raises
mean error to 31.47\% and lowers coverage to 18/23. This counterexample shows
why choosing a model by one favorable small-sample statistic would be
misleading. No winner is selected from the audit.

Six MaNGA H$\alpha$ velocity maps provide a further data-availability check.
All are processed, but none passes the specified quality/extent/flatness
requirements for the flat-H\,I-speed calibration. Distance is therefore
withheld. This outcome concerns those optical maps and those gates, not a
claim that H$\alpha$ is universally unsuitable for kinematic distances.
Neither these maps nor the three illustrative equilibrium targets supply
missing image-only distance labels for the redshift experiment.

\section{Reproducibility, outputs and interpretation}
The companion repository provides publicly available source, model metadata,
machine-readable measurements and figure provenance:
\url{https://github.com/nalin-dhiman/Galaxy-arm-project}.
Original archive products are acquired using their credited public interfaces.
The reproducibility sequence is to acquire/verify identities and native data;
validate coordinates, PSFs and coverage; apply the frozen region model;
assemble observable features by identity; apply the fixed redshift regressors
and calibration; and compute galaxy/group outcomes with the declared
denominators. A fresh scientific experiment requires a new fixed sample and
protocol rather than overwriting the reported frozen evaluation.

Per-object region outputs retain the valid mask, score map, thresholded
display map, angular sampling, coverage, feature schema and model identity.
Redshift outputs retain the predicted catalog redshift, interval endpoints,
feature-family identity, missing-input flags and support status. Accepted
science results and diagnostic candidates have separate meanings. An
out-of-range or incompatible input is reported explicitly rather than
silently converted into a physical distance.

Geometry outputs retain projected/deprojected units, supplied center and disk
geometry, candidate versus matched segment counts and the reference convention.
Physical conversions retain the supplied distance and its origin. Dynamical
outputs retain baryonic assumptions, boundary/Poisson diagnostics and the
kinematic observable definition. These records are necessary to prevent a
good image visualization, a successful numerical equilibrium or a wide covered
redshift interval from being mistaken for a validated absolute-distance
measurement.

\begingroup\small
\bibliographystyle{plainnat}
\bibliography{references}
\endgroup
\end{document}
